\subsection{Conjugacy Classes and the Class Equation}

\begin{definition} \label{def:conjugacy-class-centralizer} \hlt{Conjugacy Classes and Centralizers}\\
Let $G$ be a group.
\begin{enumerate}[label=(\roman*)]
\setlength{\itemsep}{0pt}
\item The \hlt{conjugation action} of $G$ on itself is $g\cdot x\defeq gxg^{-1}$.
\item The \hlt{conjugacy class} of $x\in G$ is $\{x\}^{G}\defeq\{gxg^{-1}:g\in G\}$.
\item The \hlt{centralizer} of $x\in G$ is $C_G(x)\defeq\{g\in G:gx=xg\}$.
\item The \hlt{center} of $G$ is $Z(G)\defeq\{g\in G:(\forall\,x\in G)\,gx=xg\}$.
\end{enumerate}
\end{definition}

\begin{theorem} \label{thm:finite-group-class-equation} \hlt{Class Equation}\\
Let $G$ be finite and let $S\subseteq G$ contain one representative of each noncentral conjugacy class. Then $|G|=|Z(G)|+\sum\limits_{x\in S}[G:C_G(x)]$.
\begin{proof}
The conjugacy classes form a partition $G=Z(G)\sqcup\bigsqcup\limits_{x\in S}\{x\}^{G}$. Indeed, $\abs{\{x\}^{G}}=1 \iff (\forall\,g\in G)\,gxg^{-1}=x \iff x\in Z(G)$.\\
For $x\in S$, $\abs{\{x\}^{G}}=[G:C_G(x)]$ by Theorem \ref{thm:orbit-stabilizer-finite-group-actions} with conjugation action from Definition \ref{def:conjugacy-class-centralizer}. Take cardinalities for formula.\\
\hlt{Structural Commentary}: $G$ decomposes into central singleton classes and noncentral conjugacy classes.\\
\hlt{Key Insight}: Count conjugacy classes by orbit-stabilizer.
\end{proof}
\end{theorem}
